% ============================================================
%  02_estado_del_arte.tex — Capítulo 2: Estado del arte
% ============================================================

\chapter{Estado del arte}
\label{cap:estado-del-arte}

\textit{[Este capítulo resume la literatura relacionada. Usa citas con \texttt{\textbackslash cite\{clave\}}
y añade las referencias a \texttt{bibliografia/referencias.bib}.]}

\section{Agricultura de precisión}
\label{sec:agricultura-precision}

\textit{[Revisión de la agricultura de precisión: sensores, IoT, automatización de invernaderos.]}
% Ejemplo de cita: \cite{smith2020}

\section{Sistema de control y monitoreo en invernaderos}
\label{sec:control-invernaderos}

\textit{[Revisión de arquitecturas de control, protocolos de comunicación y plataformas.]}

\section{Ontologías y representación del conocimiento}
\label{sec:ontologias}

\textit{[Ontologías aplicadas a dominios agrícolas, lenguajes (OWL, RDF, RDFS) y herramientas como Protégé.]}

\section{Sistemas inteligentes y ciencia de datos}
\label{sec:ia-datos}

\textit{[Modelos de aprendizaje automático aplicados a predicción, optimización y control climático.]}

\section{Análisis crítico y oportunidades}
\label{sec:analisis-critico}

\textit{[¿Qué vacíos existen? ¿Qué aporta esta tesis frente a lo revisado?]}
